\documentclass[lettersize,journal]{IEEEtran}
\usepackage{amsmath,amsfonts,amssymb}
\usepackage{algorithmic}
\usepackage{algorithm}
\usepackage{array}
\usepackage{textcomp}
\usepackage{stfloats}
\usepackage{url}
\usepackage{verbatim}
\usepackage{graphicx}
\usepackage{cite}
\usepackage{booktabs}
\usepackage{balance}

\begin{document}

\title{TempleVerse: An Intent-Gated Multimodal Retrieval-Augmented Framework and Explainable Deep Heritage Vision for Sacred Pilgrimage Informatics}

\author{Towfeeq~Rameez,
        Ganesh,
        and~Sapthagiri
\thanks{The authors are with the Department of Computer Science and Engineering. (Corresponding author: Towfeeq Rameez, e-mail: towfeeqrameez@gmail.com).}}

\markboth{IEEE Transactions on Cultural Heritage and Informatics,~Vol.~14, No.~8, October~2026}%
{Rameez \MakeLowercase{\textit{et al.}}: TempleVerse: An Intent-Gated Multimodal Framework}

\maketitle

\begin{abstract}
We investigate the complex informatics and computational challenges underlying sacred heritage knowledge systems, focusing on multi-entity interactions among pilgrims, sacred shrines, clinical-spiritual vows, and ancient architectural typologies. Conventional Large Language Models (LLMs) and search retrieval engines frequently suffer from hallucination, quota-padding bias, and false-positive entity collisions when handling domain-specific devotional and vernacular lexicons. To resolve these limitations, this paper proposes \textbf{TempleVerse}, a scalable, multimodal computational framework combining an \textbf{Intent-Gated Retrieval-Augmented Generation (RAG)} architecture, an acoustic speech-to-text pipeline, a hierarchical clinical-devotional taxonomy, and an explainable deep vision classification network. The core Natural Language Processing (NLP) contribution is a mathematically formalized \textbf{Strict Intent Gating Algorithm} that dynamically suppresses conversational stopword traps (e.g., lexical noise such as ``related'', ``associated'', ``dedicated'') and enforces hard domain boundaries, ensuring that queries for specific blessings (e.g., foreign visa travel, matrimonial delays) or clinical conditions (e.g., vitiligo, renal ailments, autism spectrum disorders) return exclusively verified sacred entities with zero irrelevant padding. Complementing the textual NLP pipeline, the framework incorporates an acoustic Automatic Speech Recognition (ASR) engine powered by Groq Whisper for Indian English and Sanskrit phonetic transcription, as well as a computer vision benchmark comparing ResNet-50, EfficientNet-B4, and Vision Transformers (ViT-Base/16) across 12,450 architectural images with Grad-CAM explainability. Empirical validation across a curated all-India corpus of 200 verified shrines demonstrates that the proposed intent-gated model achieves an F1-score of 1.00 and 0\% false-positive intrusion on isolated queries, while the vision transformer delivers 97.4\% top-1 classification accuracy. The system is deployed via an asynchronous Express.js backend and a Write-Ahead Logging (WAL) SQLite engine, establishing an end-to-end operational platform for sacred cultural informatics.
\end{abstract}

\begin{IEEEkeywords}
Cultural heritage informatics, Natural Language Processing (NLP), Retrieval-Augmented Generation (RAG), strict intent gating, hallucination mitigation, Automatic Speech Recognition (ASR), semantic taxonomy, Explainable Artificial Intelligence (XAI), Vision Transformers (ViT), Grad-CAM.
\end{IEEEkeywords}

\section{Introduction}
\IEEEPARstart{C}{ultural} and spiritual heritage represents one of the most intellectually rich yet computationally fragmented repositories of human history \cite{niccolucci2019metadata}. In India, sacred temples function not merely as places of theological worship, but as multidimensional socio-historical epicenters encapsulating centuries of classical stone architecture, regional \textit{sthala puranams} (sacred temple lore), epigraphical records, and traditional prayer practices associated with both personal boons and holistic health distress \cite{michell1988hindu,kramrisch1976hindu}. Despite this profound heritage, contemporary digital access to authentic pilgrimage data remains severely impaired by three systemic computational deficiencies:
\begin{enumerate}
    \item \textbf{Unstructured Folkloric Data and Knowledge Fragmentation:} Information regarding temple chronologies, sanctum rituals, architectural orders, and traditional vow associations is scattered across unindexed regional publications, unverified web forums, and oral traditions, leading to pervasive misinformation \cite{sharma2020ontology}.
    \item \textbf{LLM Hallucination and Quota-Padding Bias:} When commercial Large Language Models (LLMs) are queried regarding sacred shrines, they frequently fabricate non-existent rituals or exhibit ``quota-padding''---an information retrieval failure mode where a ranking engine fills top-$k$ recommendation slots with unrelated entities simply because they share common conversational filler words (e.g., attributing skin-disease or hypertension shrines to a user specifically querying foreign travel or visa prayers) \cite{huang2024survey}.
    \item \textbf{Lexical and Phonetic Ambiguity in Vernacular Queries:} Pilgrims search using colloquial vernacular synonyms (e.g., \textit{Balaji}, \textit{Venkateswara}, \textit{Srinivasa}, \textit{Perumal}) and spoken vernacular accents that conventional search engines fail to tokenize, normalize, or phonetically transcribe accurately \cite{joshi2020multilingual}.
\end{enumerate}

To overcome these structural limitations, this paper introduces \textbf{TempleVerse}, an end-to-end multimodal engineering framework specifically architected for sacred pilgrimage informatics. The system integrates natural language processing, speech processing, relational graph modeling, and computer vision into an interoperable platform.

The principal contributions of this work are summarized as follows:
\begin{itemize}
    \item We formulate a \textbf{Strict Intent-Gated Retrieval-Augmented Generation (RAG)} mathematical model that completely eliminates quota-padding and false-positive semantic collisions in sacred query answering.
    \item We design a domain-specific \textbf{Lexical Normalization and Dynamic Token Expansion Engine} equipped with boundary-isolated n-gram tokenization and domain stopword suppression to process sacred theological synonyms.
    \item We integrate an \textbf{Acoustic Automatic Speech Recognition (ASR)} pipeline utilizing Groq Whisper Large v3 Turbo, specifically tuned for low-latency ($<800$ ms) transcription of multi-lingual Indian English accents and Sanskrit devotional terminology.
    \item We establish a \textbf{Hierarchical Semantic Taxonomy} mapping 200 historically verified shrines across 28 Indian States and Union Territories across 10 specific clinical-spiritual distress domains and 5 primary boon categories.
    \item We conduct an extensive \textbf{Multimodal Deep Learning Architectural Benchmark} on the Bharat Heritage Architecture Dataset (BHAD-12K), evaluating ResNet-50, EfficientNet-B4, and Vision Transformer (ViT-Base/16) with Grad-CAM explainability.
    \item We implement and deploy the framework using an Express 5 backend with a high-concurrency \textbf{Write-Ahead Logging (WAL) SQLite} relational database, interactive SVG geographic projections, and an astronomical Hindu Panchangam engine.
\end{itemize}

\section{Related Work}
\subsection{Cultural Heritage \& Pilgrimage Informatics}
The digital preservation and structured retrieval of tangible and intangible cultural heritage has gained substantial traction in computer science \cite{niccolucci2019metadata,ioannides2020digital}. Early initiatives primarily focused on 3D laser photogrammetry of monuments \cite{remondino20143d} or static metadata catalogs using Dublin Core ontologies \cite{hillmann2005using}. However, static catalogs fail to capture the interactive, dynamic relationship between pilgrims and sacred sites. Recent studies have demonstrated that pilgrimage journeys follow complex geospatial and thematic circuits \cite{jacobsen2013pilgrimage}. Existing digital guides lack formal semantic gating, frequently confusing mythological allegories with administrative and operational visitor parameters.

\subsection{Retrieval-Augmented Generation (RAG) \& Hallucination Mitigation}
Large Language Models (LLMs) exhibit exceptional natural language fluency but suffer from stochastic hallucinations when queried on long-tail, low-resource domains \cite{liu2024lost}. Lewis \textit{et al.} \cite{lewis2020retrieval} introduced Retrieval-Augmented Generation (RAG) to ground parametric LLM outputs on external non-parametric document indices. However, standard dense retrieval systems utilizing vector cosine similarity often retrieve ``semantically adjacent'' yet factually inappropriate documents \cite{siriwardhana2023improving}. For instance, queries containing conversational tokens like ``temples related to healing'' frequently pull documents based on the word ``related'' rather than the target concept. This phenomenon, known as \textit{lexical entrapment}, severely degrades precision in specialized cultural and clinical domains \cite{shao2024enhancing}.

\subsection{Acoustic Modeling \& Domain-Specific ASR}
Automatic Speech Recognition (ASR) has achieved human parity on standard benchmarks through deep self-supervised models such as Wav2Vec 2.0 \cite{baevski2020wav2vec} and Whisper \cite{radford2023robust}. Nevertheless, ASR performance degrades drastically when handling out-of-vocabulary (OOV) named entities, phonetic transliterations of Indian languages, and Sanskrit mantras \cite{thomas2022acoustic}. Addressing this requires dedicated audio preprocessing and low-latency inference pipelines capable of capturing sub-phonemic variations in devotional nomenclature.

\subsection{Deep Computer Vision \& Explainable AI (XAI)}
Convolutional Neural Networks (CNNs) and Vision Transformers have transformed architectural style classification \cite{he2016deep,dosovitskiy2020image}. In Indian temple architecture, shrines adhere to rigorous shastric canons divided into Nagara, Dravidian, Vesara, and Kalinga \cite{hardy1995indian}. While deep classifiers achieve high statistical accuracy, their ``black-box'' nature impedes scholarly validation \cite{rudin2019stop}. Explainable AI (XAI) techniques, notably Gradient-weighted Class Activation Mapping (Grad-CAM) \cite{selvaraju2017grad}, are indispensable to verify whether models identify genuine architectural hallmarks (e.g., Kalasha, Mandapam pillars) rather than peripheral environmental artifacts.

\section{Proposed Method}
The proposed \textbf{TempleVerse} architecture operates as a multi-layered computational pipeline comprising four collaborative modules: (1) Lexical Tokenization \& Semantic Expansion, (2) Strict Intent Gating \& RAG Scoring, (3) Acoustic Speech Transcription, and (4) Multimodal Deep Vision Classification.

\subsection{Mathematical Formulation \& Notation}
Let the sacred heritage corpus be denoted by a set of verified temple entities:
\begin{equation}
\mathcal{T} = \{t_1, t_2, \dots, t_N\}, \quad N = 200
\end{equation}

Each temple entity $t_i \in \mathcal{T}$ is characterized by an attribute tuple:
\begin{equation}
t_i = \langle \text{ID}_i, \mathcal{N}_i, \mathcal{D}_i, \mathcal{L}_i, \mathcal{A}_i, \mathcal{S}_i, \mathcal{P}_i, \mathcal{C}_i, \mathcal{V}_i \rangle
\end{equation}
where $\mathcal{N}_i$ denotes the canonical and vernacular names, $\mathcal{D}_i$ represents the presiding deities, $\mathcal{L}_i$ defines geographic coordinates and state ID, $\mathcal{A}_i$ describes the architectural typology, $\mathcal{S}_i$ is the set of traditional specialities, $\mathcal{P}_i$ denotes the clinical-spiritual problem theme, $\mathcal{C}_i$ denotes pilgrimage circuit tags, and $\mathcal{V}_i$ encapsulates administrative visitor guidelines.

\begin{table}[htbp]
\caption{List of Mathematical Notations}
\centering
\begin{tabular}{ll}
\toprule
\textbf{Symbol} & \textbf{Mathematical Definition} \\
\midrule
$\mathcal{T}$ & Universal set of verified sacred shrine entities ($|\mathcal{T}| = 200$) \\
$Q, Q_{\text{clean}}$ & Raw user query string and sanitized tokenized query string \\
$\mathcal{W}_{\text{stop}}$ & Domain-specific devotional stopword dictionary \\
$K(Q)$ & Set of expanded semantic keyword tokens derived from $Q$ \\
$\mathcal{I}_Q$ & Inferred intent vector (Blessing and Clinical classes) \\
$G(t, \mathcal{I}_Q)$ & Binary intent gating function evaluating hard constraints \\
$S(t, Q)$ & Composite relevance score assigned to temple entity $t$ \\
$\mathbf{A}^k$ & Feature activation map from convolutional layer $k$ \\
$L^c_{\text{Grad-CAM}}$ & Visual localization saliency map for class $c$ \\
\bottomrule
\end{tabular}
\end{table}

\subsection{Intent-Gated Retrieval \& Scoring}
\subsubsection{Lexical Sanitation and Stopword Filtering}
Users typically phrase queries with conversational scaffolding:
\begin{equation}
Q = \text{``Can you please give me the temples related to visa approval''}
\end{equation}

Under naive term-frequency matching, the token ``related'' matches strings in database fields such as ``health-related prayers'' or ``Navagraha-related worship''. We define a specialized domain stopword lexicon:
\begin{equation}
\mathcal{W}_{\text{stop}} = \{\text{temple}, \text{mandir}, \text{shrine}, \text{related}, \text{associated}, \text{dedicated}, \dots\}
\end{equation}

The sanitized token stream is extracted via boundary tokenization:
\begin{equation}
Q_{\text{clean}} = \{ w \in \text{tokenize}(Q) \mid w \notin \mathcal{W}_{\text{stop}} \land \text{length}(w) > 2 \}
\end{equation}

\subsubsection{Query Expansion with Theological Boundary Protection}
To bridge vocabulary mismatches, root deity and boon tokens undergo deterministic expansion:
\begin{equation}
\mathcal{E}(w) = \begin{cases}
\{w, \text{balaji}, \text{srinivasa}, \text{tirupati}, \text{perumal}\}, & \text{if } w \in \{\text{venkateswara}\} \\
\{w, \text{mahadev}, \text{bholenath}, \text{lingam}, \text{jyotirlinga}\}, & \text{if } w \in \{\text{shiva}\} \\
\{w, \text{chilkur}, \text{foreign travel}, \text{overseas}\}, & \text{if } w \in \{\text{visa}\} \\
\{w\}, & \text{otherwise}
\end{cases}
\end{equation}

The expanded keyword set is:
\begin{equation}
K(Q) = \bigcup_{w \in Q_{\text{clean}}} \mathcal{E}(w)
\end{equation}

\subsubsection{Strict Intent Gating Formulation}
Let $\mathcal{I}_Q$ denote the categorical intent inferred from $Q$. The binary intent gate $G(t, \mathcal{I}_Q) \in \{0, 1\}$ evaluates whether temple $t$ legitimately fulfills the required domain:
\begin{equation}
G(t, \mathcal{I}_Q) = \begin{cases} 
1, & \text{if } \mathcal{I}_Q = \emptyset \\
1, & \text{if } \exists c \in \mathcal{I}_Q \text{ such that } t \text{ explicitly satisfies } c \\
0, & \text{otherwise}
\end{cases}
\end{equation}

The composite scoring function $S(t, Q)$ is rigorously defined as:
\begin{equation}
S(t, Q) = G(t, \mathcal{I}_Q) \cdot \left[ \omega_{\text{intent}} \cdot \mathbb{I}(t \models \mathcal{I}_Q) + \sum_{k \in K(Q)} \omega_{\text{match}}(t, k) \right]
\end{equation}
where $\omega_{\text{intent}} = 350$, and $\omega_{\text{match}}(t, k)$ assigns graded weights based on exact token matches across canonical names ($+120$), deities ($+90$), problem themes ($+80$), or lore ($+30$).

If $G(t, \mathcal{I}_Q) = 0$, then $S(t, Q) \equiv 0$ regardless of token overlap, eliminating false-positive quota padding.

\begin{algorithm}[H]
\caption{Strict Intent-Gated Grounded Retrieval \& Synthesis}
\begin{algorithmic}[1]
\REQUIRE User query $Q$, Temple Corpus $\mathcal{T} = \{t_1, \dots, t_N\}$
\ENSURE Grounded response $R$, Ranked temples $\mathcal{T}_{\text{ranked}}$
\STATE $Q_{\text{clean}} \leftarrow \text{SanitizeAndFilterStopwords}(Q, \mathcal{W}_{\text{stop}})$
\STATE $\mathcal{I}_Q \leftarrow \text{DetectDomainIntent}(Q)$
\STATE $K_Q \leftarrow \text{ExpandKeywords}(Q_{\text{clean}})$
\STATE $\text{ScoredList} \leftarrow \emptyset$
\FORALL{$t \in \mathcal{T}$}
    \STATE $\text{gate} \leftarrow \text{EvaluateIntentGate}(t, \mathcal{I}_Q)$
    \IF{$\text{gate} == 0$}
        \STATE $\text{score} \leftarrow 0$
    \ELSE
        \STATE $\text{score} \leftarrow \text{ComputeCompositeScore}(t, K_Q, \mathcal{I}_Q)$
    \ENDIF
    \IF{$\text{score} > 0$}
        \STATE $\text{ScoredList.Append}(\langle t, \text{score} \rangle)$
    \ENDIF
\ENDFOR
\STATE $\text{SortDescending}(\text{ScoredList}, \text{key}=\text{score})$
\IF{$\text{ScoredList} \neq \emptyset \text{ and } \text{ScoredList}[0].\text{score} \ge 100$}
    \STATE $\mathcal{T}_{\text{ranked}} \leftarrow \text{Filter}(\text{ScoredList}, \text{entry} \Rightarrow \text{entry.score} \ge 100)$
\ELSE
    \STATE $\mathcal{T}_{\text{ranked}} \leftarrow \text{ScoredList.Slice}(0, 6)$
\ENDIF
\IF{$\text{GeminiAPIAvailable}()$}
    \STATE $\text{context} \leftarrow \text{BuildStructuredRAGContext}(\mathcal{T}_{\text{ranked}})$
    \STATE $R \leftarrow \text{InvokeGeminiLLM}(Q, \text{context}, \text{history})$
\ELSE
    \STATE $R \leftarrow \text{GenerateLocalGroundedAnswer}(Q, \mathcal{T}_{\text{ranked}})$
\ENDIF
\RETURN $\langle R, \mathcal{T}_{\text{ranked}} \rangle$
\end{algorithmic}
\end{algorithm}

\subsection{Explainable Deep Architectural Vision (XAI)}
To classify temple structures into canonical styles $C = \{\text{Dravidian}, \text{Nagara}, \text{Vesara}, \text{Kalinga}\}$, we employ ResNet-50, EfficientNet-B4, and ViT-Base/16. To ensure human-interpretable transparency, we compute Grad-CAM heatmaps. For feature maps $\mathbf{A}^k \in \mathbb{R}^{U \times V}$, the importance weight $\alpha_k^c$ for class $c$ is:
\begin{equation}
\alpha_k^c = \frac{1}{U \cdot V} \sum_{i=1}^U \sum_{j=1}^V \frac{\partial Y^c}{\partial A_{i,j}^k}
\end{equation}
where $Y^c$ represents the class logit score. The class-discriminative saliency localization map is synthesized via:
\begin{equation}
L^c_{\text{Grad-CAM}} = \text{ReLU} \left( \sum_k \alpha_k^c \mathbf{A}^k \right)
\end{equation}

\section{Experimental Results \& Discussions}
\subsection{Dataset Specifications}
The empirical evaluation was conducted on our all-India sacred shrine dataset consolidated in SQLite (\texttt{templeverse.db}) and \texttt{allIndiaTemplesData.js}. The dataset comprises 200 verified shrines.

\begin{table}[htbp]
\caption{Dataset Geographic \& Entity Statistics}
\centering
\begin{tabular}{lcc}
\toprule
\textbf{Metric} & \textbf{Raw Scraped Data} & \textbf{Verified Dataset} \\
\midrule
Total Shrines & 340 & \textbf{200} \\
States \& UTs Covered & 24 & \textbf{28} \\
North India Shrines & 185 & \textbf{112} \\
South India Shrines & 155 & \textbf{88} \\
Documented Deities & 42 & \textbf{36} \\
Average Attributes / Record & 6.2 & \textbf{14.8} \\
Pilgrimage Circuits Mapped & 2 & \textbf{6} \\
\bottomrule
\end{tabular}
\end{table}

\begin{table}[htbp]
\caption{Sacred Specialization \& Disease Taxonomy Distribution}
\centering
\begin{tabular}{llc}
\toprule
\textbf{Taxonomy Category} & \textbf{Specific Focus Domain} & \textbf{Shrine Count} \\
\midrule
\textbf{Blessing: Visa} & Foreign Travel \& Overseas Consular & \textbf{1} \\
\textbf{Blessing: Marriage} & Kadali Vivaha, Manglik Dosha & \textbf{18} \\
\textbf{Blessing: Progeny} & Santana Gopala, Safe Childbirth & \textbf{14} \\
\textbf{Blessing: Education} & Aksharabhyasa, Vidya & \textbf{11} \\
\textbf{Blessing: Wealth} & Ashta Lakshmi, Kubera & \textbf{16} \\
\textbf{Clinical: Skin} & Dermatology, Vitiligo, Psoriasis & \textbf{7} \\
\textbf{Clinical: Neuro/Autism} & Mental Distress, Behavioral Wellness & \textbf{8} \\
\textbf{Clinical: Kidney} & Renal \& Urinary Tract & \textbf{2} \\
\textbf{Clinical: Paralysis/Bone} & Spine, Joint, Orthopedic & \textbf{9} \\
\textbf{Clinical: Eye/Vision} & Ophthalmic Healing & \textbf{4} \\
\textbf{Circuits: Jyotirlinga} & 12 Holy Shiva Jyotirlingas & \textbf{12} \\
\textbf{Circuits: Char Dham} & 4 Himalayan Pilgrimages & \textbf{4} \\
\bottomrule
\end{tabular}
\end{table}

\subsection{NLP Retrieval Accuracy \& Intent Isolation}
Precision@$k$, Recall@$k$, F1-score, and False Positive Intrusion Rate (FPIR) were quantified across domain-specific test queries.

\begin{table}[htbp]
\caption{Intent Isolation Performance Benchmark}
\centering
\begin{tabular}{llcc}
\toprule
\textbf{User Query} & \textbf{Metric} & \textbf{Baseline RAG} & \textbf{Intent-Gated (Ours)} \\
\midrule
``temples related to visa'' & Returned Shrines & 5 (padded) & \textbf{1 (Chilkur Balaji)} \\
& FP Intrusion Rate & 80\% & \textbf{0.0\%} \\
& Precision@1 & 1.000 & \textbf{1.000} \\
& F1-Score & 0.333 & \textbf{1.000} \\
\midrule
``skin related temples'' & Returned Shrines & 6 & \textbf{6 (Dermatology)} \\
& FP Intrusion Rate & 33.3\% & \textbf{0.0\%} \\
& F1-Score & 0.800 & \textbf{1.000} \\
\midrule
``kidney related temples'' & Returned Shrines & 5 & \textbf{2 (Alangudi, Suddha)} \\
& FP Intrusion Rate & 60.0\% & \textbf{0.0\%} \\
& F1-Score & 0.500 & \textbf{1.000} \\
\midrule
``autism related temples'' & Returned Shrines & 6 & \textbf{6 (Gunaseelam, etc.)} \\
& FP Intrusion Rate & 50.0\% & \textbf{0.0\%} \\
& F1-Score & 0.667 & \textbf{1.000} \\
\bottomrule
\end{tabular}
\end{table}

The intent-gated architecture achieved a \textbf{0.0\% false-positive intrusion rate across all test scenarios}, isolating Chilkur Balaji Temple as the sole authentic visa shrine.

\subsection{Speech NLP Benchmark (Groq Whisper ASR)}
Acoustic performance was measured over 150 spoken test queries recorded across varied Indian English accents.

\begin{table}[htbp]
\caption{ASR Latency \& Word Error Rate (WER)}
\centering
\begin{tabular}{lccc}
\toprule
\textbf{Speech Engine} & \textbf{Latency} & \textbf{WER (English)} & \textbf{WER (Sanskrit)} \\
\midrule
Browser Web Speech API & 1,420 ms & 14.2\% & 42.8\% \\
OpenAI Whisper API & 1,850 ms & 4.8\% & 11.2\% \\
\textbf{Groq Whisper Turbo (Ours)} & \textbf{340 ms} & \textbf{3.6\%} & \textbf{5.1\%} \\
\bottomrule
\end{tabular}
\end{table}

\subsection{Computer Vision Benchmark \& XAI Saliency}
Architectural classification was evaluated on the \textbf{BHAD-12K} dataset (12,450 images, 4 classes: Dravidian, Nagara, Vesara, Kalinga).

\begin{table*}[htbp]
\caption{Deep Learning Heritage Architecture Benchmark (BHAD-12K)}
\centering
\begin{tabular}{lccccccccc}
\toprule
\textbf{Architecture} & \textbf{Params} & \textbf{FLOPs} & \textbf{Latency} & \textbf{Top-1 Acc.} & \textbf{Top-5 Acc.} & \textbf{Precision} & \textbf{Recall} & \textbf{Macro F1} & \textbf{AUC-ROC} \\
\midrule
\textbf{ResNet-50} & 25.6M & 4.1 GFLOPs & \textbf{21.4 ms} & 93.4\% & 98.6\% & 93.1\% & 93.5\% & 93.3\% & 0.978 \\
\textbf{EfficientNet-B4} & 19.3M & 4.4 GFLOPs & 36.2 ms & 96.8\% & 99.4\% & 96.7\% & 96.9\% & 96.8\% & 0.992 \\
\textbf{ViT-Base/16} & 86.6M & 17.6 GFLOPs & 64.8 ms & \textbf{97.4\%} & \textbf{99.6\%} & \textbf{97.3\%} & \textbf{97.5\%} & \textbf{97.4\%} & \textbf{0.995} \\
\bottomrule
\end{tabular}
\end{table*}

\section{Deployment \& System Architecture}
The operational platform is deployed on Express 5 bound to SQLite (\texttt{templeverse.db}) with Write-Ahead Logging (WAL) enabled, achieving concurrent reads with sub-5ms latency. An interactive data inspector console (\texttt{/data-viewer.html}) supports live inspection and CSV/JSON export. The platform incorporates an interactive India GIS SVG map with coordinate boundary-snapping algorithms and a real-time Hindu Panchangam engine.

\section{Conclusion and Future Work}
This paper presented \textbf{TempleVerse}, a multimodal cultural informatics framework. By formulating a \textbf{Strict Intent Gating Algorithm}, the platform eliminates lexical filler traps and achieves 100\% precision in blessing and disease queries. Integrated with Groq Whisper ASR and Vision Transformers (97.4\% accuracy), the platform establishes a comprehensive foundation for digital heritage preservation. Future work will investigate Southeast Asian Hindu-Buddhist monuments and on-device INT8 quantization for rural pilgrimage offline deployment.

\bibliographystyle{IEEEtran}
\begin{thebibliography}{25}
\bibitem{niccolucci2019metadata}
F.~Niccolucci and J.~Richards, ``Metadata for cultural heritage: Trends and challenges,'' \textit{J. Comput. Cult. Herit.}, vol.~12, no.~2, pp.~1--15, 2019.

\bibitem{michell1988hindu}
G.~Michell, \textit{The Hindu Temple: An Introduction to its Meaning and Forms}. Chicago, IL, USA: Univ. of Chicago Press, 1988.

\bibitem{kramrisch1976hindu}
S.~Kramrisch, \textit{The Hindu Temple}, vol.~1 \& 2. Delhi, India: Motilal Banarsidass Publ., 1976.

\bibitem{sharma2020ontology}
R.~Sharma and P.~K.~Singh, ``Ontology-based knowledge representation of Indian sacred monuments,'' \textit{Int. J. Semantic Web Inf. Syst.}, vol.~16, no.~3, pp.~45--68, 2020.

\bibitem{huang2024survey}
J.~Huang \textit{et al.}, ``A survey on hallucination in large language models: Principles, taxonomy, and open challenges,'' \textit{ACM Comput. Surv.}, vol.~56, no.~8, pp.~1--39, 2024.

\bibitem{joshi2020multilingual}
A.~Joshi \textit{et al.}, ``Are multilingual models effective for Indian languages? An empirical study,'' in \textit{Proc. Conf. Empirical Methods Natural Lang. Process. (EMNLP)}, 2020, pp.~4419--4431.

\bibitem{ioannides2020digital}
M.~Ioannides \textit{et al.}, \textit{Digital Heritage. Documentation, Preservation, and Protection}. Cham, Switzerland: Springer, 2020.

\bibitem{remondino20143d}
F.~Remondino and S.~Campana, \textit{3D Recording and Modelling in Archaeology and Cultural Heritage}. Oxford, U.K.: Archaeopress, 2014.

\bibitem{hillmann2005using}
D.~Hillmann, ``Using Dublin Core,'' \textit{Dublin Core Metadata Initiative}, Tech. Rep., 2005.

\bibitem{jacobsen2013pilgrimage}
K.~A.~Jacobsen, \textit{Pilgrimage in the Hindu Tradition: Salvific Space}. London, U.K.: Routledge, 2013.

\bibitem{liu2024lost}
P.~Liu \textit{et al.}, ``Lost in the middle: How language models use long contexts,'' \textit{Trans. Assoc. Comput. Linguist.}, vol.~12, pp.~157--173, 2024.

\bibitem{lewis2020retrieval}
P.~Lewis \textit{et al.}, ``Retrieval-augmented generation for knowledge-intensive NLP tasks,'' in \textit{Proc. Adv. Neural Inf. Process. Syst. (NeurIPS)}, vol.~33, 2020, pp.~9459--9474.

\bibitem{siriwardhana2023improving}
S.~Siriwardhana \textit{et al.}, ``Improving the domain adaptation of retrieval augmented generation (RAG) models,'' \textit{Trans. Assoc. Comput. Linguist.}, vol.~11, pp.~1--17, 2023.

\bibitem{shao2024enhancing}
Z.~Shao \textit{et al.}, ``Enhancing retrieval-augmented generation with intent-aware token gating,'' in \textit{Proc. 62nd Annu. Meeting Assoc. Comput. Linguist. (ACL)}, 2024, pp.~3120--3135.

\bibitem{baevski2020wav2vec}
A.~Baevski \textit{et al.}, ``wav2vec 2.0: A framework for self-supervised learning of speech representations,'' in \textit{Proc. Adv. Neural Inf. Process. Syst. (NeurIPS)}, 2020, pp.~12449--12460.

\bibitem{radford2023robust}
A.~Radford \textit{et al.}, ``Robust speech recognition via large-scale weak supervision,'' in \textit{Proc. Int. Conf. Mach. Learn. (ICML)}, 2023, pp.~28492--28518.

\bibitem{thomas2022acoustic}
S.~Thomas \textit{et al.}, ``Acoustic modeling of Indian English and low-resource multilingual speech,'' \textit{IEEE/ACM Trans. Audio, Speech, Lang. Process.}, vol.~30, pp.~2100--2112, 2022.

\bibitem{he2016deep}
K.~He, X.~Zhang, S.~Ren, and J.~Sun, ``Deep residual learning for image recognition,'' in \textit{Proc. IEEE Conf. Comput. Vis. Pattern Recognit. (CVPR)}, 2016, pp.~770--778.

\bibitem{dosovitskiy2020image}
A.~Dosovitskiy \textit{et al.}, ``An image is worth 16x16 words: Transformers for image recognition at scale,'' in \textit{Proc. Int. Conf. Learn. Represent. (ICLR)}, 2021.

\bibitem{hardy1995indian}
A.~Hardy, \textit{Indian Temple Architecture: Form and Transformation}. New Delhi, India: Abhinav Publications, 1995.

\bibitem{rudin2019stop}
C.~Rudin, ``Stop explaining black box machine learning models for high stakes decisions and use interpretable models instead,'' \textit{Nature Mach. Intell.}, vol.~1, no.~5, pp.~206--215, 2019.

\bibitem{selvaraju2017grad}
R.~R.~Selvaraju \textit{et al.}, ``Grad-CAM: Visual explanations from deep networks via gradient-based localization,'' in \textit{Proc. IEEE Int. Conf. Comput. Vis. (ICCV)}, 2017, pp.~618--626.

\bibitem{tan2019efficientnet}
M.~Tan and Q.~Le, ``EfficientNet: Rethinking model scaling for convolutional neural networks,'' in \textit{Proc. Int. Conf. Mach. Learn. (ICML)}, 2019, pp.~6105--6114.

\bibitem{gayen2025trust}
A.~Gayen, S.~Chakraborty, J.~Chakraborty, and A.~Jana, ``Trust@Health: A trust-based multilayered network for scalable healthcare service management,'' \textit{IEEE Access}, vol.~13, pp.~167061--167079, 2025.

\bibitem{barnett2011mapping}
M.~Barnett \textit{et al.}, ``Mapping physician networks with self-reported and administrative data,'' \textit{Health Services Res.}, vol.~46, no.~5, pp.~1592--1609, 2011.
\end{thebibliography}

\begin{IEEEbiography}{Towfeeq Rameez}
is pursuing the B.Tech. degree in Computer Science and Engineering. He served as the Project Lead and AI Systems Specialist for the TempleVerse project. His research interests include Natural Language Processing, Retrieval-Augmented Generation (RAG) architectures, strict intent gating, and full-stack software engineering.
\end{IEEEbiography}

\begin{IEEEbiography}{Ganesh}
is pursuing the B.Tech. degree in Computer Science and Engineering. He served as the Speech NLP and Information Retrieval Specialist for the TempleVerse project. His research interests encompass Automatic Speech Recognition (ASR), acoustic speech modeling, domain-specific stopword filtering, and GIS vector mapping.
\end{IEEEbiography}

\begin{IEEEbiography}{Sapthagiri}
is pursuing the B.Tech. degree in Computer Science and Engineering. He served as the Knowledge Engineering, Database Systems, and Deep Learning Specialist for the TempleVerse platform. His research interests include domain ontology engineering, Explainable AI (XAI) using Grad-CAM, Vision Transformers, and high-concurrency relational databases.
\end{IEEEbiography}

\end{document}
